% To be inserted in the Discussion section

\subsection{External Generalization and Remaining Limits}

The external analyses strengthen the bounded interpretation of the stochastic action model. In two independent visual categorization studies, lookalike objects elicited higher fitted $\rho$ than ordinary objects, consistent with $\rho$ indexing competition from an unchosen perceptual category. In the original German semantic dataset, SWOW-DE network proximity independently converged with the primary semantic-margin result: items that were less selectively connected to the target category relative to the competitor category had higher $\rho$.

These results should be interpreted as convergent external support, not as a replication or as causal proof. The Koenig-Robert archive did not distribute scalar face-likeness or animal-likeness ratings in the downloaded mouse-tracking files, so the external mouse-tracking result rests on the experimental lookalike-object manipulation and a binary ambiguity proxy. SWOW-DE direct cue-response coverage was sparse; the strongest SWOW result came from graph proximity, with embedding evidence in the same direction but lower coverage. Schröder et al.'s German norms were useful as lexical-semantic controls but did not contain the exact target--competitor category contrasts required to replace the primary semantic margin.

The severe benchmark remains essential. Flexible descriptive baselines and trial-fitted upper bounds can outperform constrained action models in raw path likelihood. The defensible contribution is therefore not raw curve-fitting dominance. Instead, the action model separates raw predictive flexibility from mechanistic interpretability by yielding a latent response-competition parameter, $\rho$, that tracks semantic and perceptual competition across datasets.
